\section{Architecture and Pipeline Overview}\label{sec:overview}

BaiZe is a multimodal EDA agent assembled in three stages on top of the BaiZe language backbone~\cite{baiZe-llm}: (i) from-scratch pre-training of a custom vision encoder, (ii) alignment of the encoder and the backbone into a multimodal large language model, and (iii) semiconductor-domain adaptation via reinforcement learning.
% 中文：BaiZe 是一个构建在 BaiZe 语言骨干 \cite{baiZe-llm} 之上的、由三个阶段组装的多模态 EDA 智能体：(i) 从零预训练自研视觉编码器；(ii) 将编码器与骨干对齐为多模态大语言模型；(iii) 通过强化学习完成半导体领域适配。

This section fixes the shared scaffolding---the vision encoder, the alignment projector, and the global training environment---before each stage is detailed in its own section.
% 中文：本节先固定共享脚手架——视觉编码器、对齐投影层与全局训练环境——随后各阶段再在各自章节中详述。

The language backbone is the BaiZe $\approx$2.2B hybrid model, fixed and imported from its companion paper~\cite{baiZe-llm}; we do not re-train or re-describe it here, and keep it frozen throughout every stage.
% 中文：语言骨干是 BaiZe 约 22 亿参数的混合模型，从其姊妹论文 \cite{baiZe-llm} 固定导入；本文不重新训练也不赘述它，并在所有阶段中保持冻结。

\subsection{Vision Encoder}\label{sec:backbone-enc}
The vision encoder is pre-trained from scratch in stage~(i) and is selected from iso-parameter candidates at the $500$--$600$\,M scale; Section~\ref{sec:vision} reports the comparison and its outcome.
% 中文：视觉编码器。视觉编码器在第 (i) 阶段从零预训练，是在 $500$–$600$M 参数量级的同参数候选模型中选出的；第 \ref{sec:vision} 节报告该对比及其结果。

\subsection{Pipeline and Global Environment}\label{sec:pipeline}
Table~\ref{tab:pipeline} summarises the three-stage pipeline.
% 中文：流水线与全局环境。表 \ref{tab:pipeline} 汇总了三阶段流水线。
The vision encoder and alignment stages run on a single $8\times$H100 node (CUDA~12.8, PyTorch~2.8.0, Python~3.10); the aligned multimodal model uses a $196$-token-per-image budget throughout.
% 中文：视觉编码器与对齐阶段运行在单 $8\times$H100 节点上（CUDA 12.8、PyTorch 2.8.0、Python 3.10）；对齐后的多模态模型全程使用每图 196 token 的预算。

\begin{table}[t]
\centering
\caption{BaiZe multimodal pipeline overview.}
\label{tab:pipeline}
\begin{tabular}{llll}
\toprule
\textbf{Stage} & \textbf{Task} & \textbf{Status} & \textbf{Section} \\
\midrule
(i) & Vision encoder & complete & \ref{sec:vision} \\
(ii) & MLLM alignment & in progress & \ref{sec:mllm} \\
(iii) & Domain adapt.\ (RL) & design & \ref{sec:domain} \\
\bottomrule
\end{tabular}
\end{table}